% === Begin Label Data ===
% ID: 878
% 难度星级: 
% 标签: 
% 备注: 新定义
% 组卷引用次数: 0
% === End  Label Data ===

\begin{problem}{2022}{G}{北京卷}{21}{数列}
已知 $Q: a_1, a_2, \dots, a_k$ 为有穷整数数列. 给定正整数 $m$, 若对任意的 $n \in \{1, 2, \dots, m\}$, 在 $Q$ 中存在 $a_i, a_{i+1}, a_{i+2}, \dots, a_{i+j} (j \geqslant 0)$, 使得 $a_i + a_{i+1} + a_{i+2} + \cdots + a_{i+j} = n$, 则称 $Q$ 为 $m-$连续可表数列.

（I）判断 $Q: 2, 1, 4$ 是否为 $5-$连续可表数列? 是否为 $6-$连续可表数列? 说明理由;

（II）若 $Q: a_1, a_2, \dots, a_k$ 为 $8-$连续可表数列, 求证: $k$ 的最小值为 $4$;

（III）若 $Q: a_1, a_2, \dots, a_k$ 为 $20-$连续可表数列, 且 $a_1 + a_2 + \cdots + a_k < 20$, 求证: $k \geqslant 7$.
\end{problem}

\begin{answer}
（I）是 $5-$连续可表数列，不是 $6-$连续可表数列；（II）证明见解析；（III）证明见解析。
\end{answer}

\begin{solutions}
（I）若 $m=5$, 则对于任意的 $n \in \{1, 2, 3, 4, 5\}$,
$a_2 = 1, a_1 = 2, a_1 + a_2 = 2 + 1 = 3, a_3 = 4, a_2 + a_3 = 1 + 4 = 5$,
所以 $Q$ 是 $5-$连续可表数列;
由于不存在任意连续若干项之和相加为 $6$,
所以 $Q$ 不是 $6-$连续可表数列;

（II）假设 $k$ 的值为 $3$, 则 $a_1, a_2, a_3$ 最多能表示
$a_1, a_2, a_3, a_1+a_2, a_2+a_3, a_1+a_2+a_3$, 共 $6$ 个数字,
与 $Q$ 是 $8-$连续可表数列矛盾, 故 $k \geqslant 4$;
现构造 $Q: 4, 2, 1, 5$ 可以表达出 $1, 2, 3, 4, 5, 6, 7, 8$ 这 $8$ 个数字,
即存在 $k=4$ 满足题意.
故 $k$ 的最小值为 $4$.

（III）当 $k=7$ 时, 数列 $1, 2, 4, 5, 8, -2, -1$ 满足题意,
当 $k>7$ 时, 数列 $1, 2, 4, 5, 8, -2, -1, a_k = \cdots = a_n = 0$, 所以 $k \geqslant 7$ 符合题意,
从 $5$ 个正整数中, 取一个数字只能表示自身, 最多可表示 $5$ 个数字,
取连续两个数字最多能表示 $4$ 个数字, 取连续三个数字最多能表示 $3$ 个数字,
取连续四个数字最多能表示 $2$ 个数字, 取连续五个数字最多能表示 $1$ 个数字,
所以对任意给定的 $5$ 个整数, 最多可以表示 $5+4+3+2+1=15$ 个正整数, 不能表示 $20$ 个正整数, 即 $k \geqslant 6$.

若 $k=6$, 最多可以表示 $6+5+4+3+2+1=21$ 个正整数,
下证明 $k=6$ 不成立即可.
由于 $Q$ 为 $20-$连续可表数列, 且 $a_1+a_2+\cdots+a_k < 20$,
所以其中必有且只有一项为负数, 而且其他 $5$ 项不重复, 其他项的连续可表项也不重复.
不妨设该项为 $a_i (i \in \{1, 2, 3, 4, 5, 6\})$,
因此数列中一定存在若干项正数的和为 $20$.
根据对称性我们只需考察 $i=1, 2, 3$ 的情况.
(i) 当 $i=2$ 时, 由于 $a_1+a_2+\cdots+a_6 < 20$, 而数列中一定有若干个连续整数和为 $20$, 因此 $a_1$ 或 $a_3+a_4+a_5+a_6$ 为连续若干个数的和中最大的数.
因此除了 $a_i$ 以外, 其它连续若干个数的和分别表示 $1, 2, \cdots, 20$.
故我们有 $a_1+a_2, a_2+a_3+a_4+a_5+a_6 \in \{1, 2, \cdots, 20\}$.
若 $a_1$ 为连续若干个数的和中最大的数, 则 $a_1+a_2+\cdots+a_6 > a_1$, 矛盾.
若 $a_3+a_4+a_5+a_6$ 为连续若干个数的和中最大的数, 同理可得矛盾;
(ii) 当 $i=3$ 时, 与 $i=2$ 同理, 不符合题意;
(iii) 当 $i=1$ 时, 则连续若干个数的和中最大的数为 $a_2+a_3+\cdots+a_6=20$.

下面有两种角度, 方法一考虑剩下的数:
有 $19 \in \{a_1+a_2+\cdots+a_6', a_2'+a_3+a_4+a_5', a_3+a_4+a_5+a_6\}$.
而由前分析可知 $a_1+a_2 > 0$, 因此 $19 \in \{a_1+a_2+\cdots+a_6', a_2+a_3+a_4+a_5\}$.
\circled{1} 若 $a_1+a_2+\cdots+a_6=19$, 则 $a_1=-1$.
*如果 $a_2+a_3+a_4+a_5=18$, 那么 $a_6=2$, 并且此时 $a_3+a_4+a_5+a_6=17$,
有 $a_2=3$, 则 $a_6=a_1+a_2$, 矛盾;
**如果 $a_3+a_4+a_5+a_6=18$, 那么 $a_2=2$, 并且此时 $a_2+a_3+a_4+a_5=17$, 有 $a_6=3$.
下面我们考察 $a_3, a_4, a_5$.
注意到 $a_3+a_4+a_5=15$ 且 $a_3, a_4, a_5 \neq 1, 2, 3$,
则 $\{a_3, a_4, a_5\} = \{4, 5, 6\}$, 并且由于 $a_2+a_3=a_3+2$ 且 $a_1+a_2+a_3=a_3+1$,
由不重复的原则可知 $a_3=6$. 如有 $a_4=4, a_5=5$,
那么 $a_5+a_6=8=a_2+a_3$, 矛盾; 如有 $a_4=5, a_5=4$,
那么 $a_6+a_6=7=a_1+a_2+a_3$, 矛盾. 故该假设不成立;
\circled{2} 若 $a_2+a_3+a_4+a_5=19$, 那么 $a_6=1$, 并且有 $a_1+a_2+\cdots+a_6=18$
并且 $a_3+a_4+a_5+a_6=17$, 则此时 $a_1=-2, a_2=3$,
那么此时 $a_1+a_2=a_6$, 矛盾, 故该假设不成立.
因此, $k \neq 6$.
综上所述 $k \geqslant 7$.

结合 $a_2+a_3+\cdots+a_6=20$, 由于 $a_1+\cdots+a_6 < 19$
若 $a_1+a_2+a_3+a_4=19, a_5=1$, 则 $a_5+a_6 < 0$, 矛盾.
若 $a_2+a_3+a_4+a_5=19, a_1=1$, 由于 $a_6 < -2, a_1+\cdots+a_6 < 18$
只可能有 $a_1+a_2+a_3+a_4=18, a_5=2$, 此时 $a_5+a_6 < 0$, 矛盾.
则 $a_6 < -2$ 时, 没有满足要求的数列.
\circled{2} 若 $a_6=-2$
若存在 $a_i=2$, 若 $i=5$, 则 $a_5+a_6=0$,
若 $i \neq 5$ 此时 $a_1+\cdots+a_6=a_{i+1}+\cdots+a_5$ 矛盾.
数列中不存在一项为 $2$, 必存在连续两项和为 $2$.
则只可能有 $a_5=4$,
考虑 $19$ 的来源, 只可能是 $a_2+a_3+a_4+a_5=19$,
则 $a_1=1$,
而 $a_2+a_3+a_4=15$, 且这些项不可能是 $1, 2, 4$.
平均值原理可得至少有一项是 $3$,
若是第二项, 则 $a_1+a_2=a_5$, 矛盾;
若是第三项, 考虑 $5$ 的来源, 矛盾;
若是第四项, 考虑 $6$ 的来源, 矛盾;
则 $a_6=-2$ 时, 没有满足要求的数列.
\circled{3} 若 $a_6=-1$
同理可证, 数列中不存在一项为 $1$, 必存在连续两项和为 $1$.
考虑 $2$ 的来源 $a_5=2$.
则 $a_1+a_2+a_3+a_4=18$
剩余四项必为 $3, 4, 5, 6$.
若 $a_4 \neq 6$, 则 $a_4+a_5+a_6=a_4+1=4$ 或 $5$ 或 $6$, 矛盾;
故 $a_4=6, a_4+a_5+a_6=7$, 故 $3, 4$ 不相邻, 数列有以下 $2$ 种情况:
i. $3, 5, 4, 6, 2, -1$: 对于这种情况 $a_1+a_2=a_4+a_5$, 矛盾;
ii. $4, 5, 3, 6, 2, -1$: 对于这种情况 $a_2+a_3=a_4+a_5$, 矛盾;
则 $a_6=-1$ 时, 没有满足要求的数列.
综上, Q.E.D.
\end{solutions}